\documentclass[10pt,journal]{IEEEtran}
\usepackage{graphicx}\usepackage{booktabs}\usepackage{url}
\title{ArchGuard-APR: Architecture-Constraint-Aware Multi-Agent Program Repair Using Repository Graphs and Execution Feedback}
\author{Majid Feiz\\Affiliation, email, and ORCID: HUMAN\_CONFIRMATION\_REQUIRED}
\begin{document}\maketitle
\begin{abstract}
Repository-level repair must preserve structural as well as behavioral constraints. ArchGuard-APR
integrates deterministic repository analysis, immutable graphs, explicit architecture profiles,
graph-guided context, role-based orchestration, patch validation, isolated execution feedback, cost
controls, and canonical evidence. In an exploratory paired evaluation, six executions were planned,
four completed, and two attrited before provider use. Baseline and ArchGuard each achieved 0/2
successful repairs; all four completed outputs were invalid patches. Usage totaled 35,431 input and
2,737 output tokens at USD 0.1296325. No effectiveness advantage was observed. Patch validity and
infrastructure robustness are the main future-work targets.
\end{abstract}
\begin{IEEEkeywords}automated program repair, software architecture, repository graph, agents, reproducibility\end{IEEEkeywords}

\section{Introduction}
Repository-level automated program repair (APR) requires context selection, admissible patch
generation, and controlled validation. A test-passing candidate may still violate architecture or
security policy. ArchGuard-APR makes these boundaries explicit and fail-closed. Its contribution is an
implemented architecture-aware prototype and reproducible research infrastructure; its scientific
evidence is a small negative and inconclusive feasibility study, not an effectiveness demonstration.

\section{Background and Related Work}
Tool-using LLM repair agents have been explored by RepairAgent~\cite{Bouzenia2024RepairAgent}.
AutoCodeRover uses structured repository search~\cite{Zhang2024AutoCodeRover}, while SpecRover adds
intent extraction and review~\cite{Ruan2024SpecRover}. These primary records motivate agents,
repository context, and review, but do not establish ArchGuard-APR's effectiveness or priority. Search
leads lacking independently verified metadata are excluded.

\section{Research Questions}
RQ1 asks whether architecture constraints, repository graphs, retrieval, and execution feedback can be
integrated reproducibly. RQ2 asks what operational limitations emerge on qualified defects. RQ3 asks
what outcomes, tokens, and costs are observed. No question asserts superiority.

\section{Design and Repository Representation}
Entrypoints call application services over a vendor-neutral domain. Infrastructure adapters implement
Git, artifacts, containers, and providers. Analysis parses source without importing target code and
emits a canonical snapshot. Graph construction consumes that snapshot without rereading source.
Modules, definitions, containment, and imports support deterministic queries and context retrieval.
The representation is primarily static and cannot model all reflective or runtime behavior.

\section{Architecture Constraints and Repair Workflow}
Versioned profiles map modules to human-reviewed components and constrain dependencies. Unassigned or
ambiguous modules invalidate evaluation. The workflow coordinates analysis, retrieval, provider,
validation, and evaluator roles through typed states. Provider calls require frozen configuration,
acknowledgements, and cost guards. Patch parsing and security policy precede isolated evaluation.
Invalid candidates remain terminal invalid-patch records rather than being silently normalized.

\section{Context, Execution, Security, and Reproducibility}
Bounded graph queries select module and symbol neighborhoods and emit manifest fingerprints.
Execution uses resource, path, network, and output controls. Isolation reduces but does not eliminate
risk. Canonical JSON, stable ordering, call accounting, and source-free diagnostics preserve evidence
without publishing target source, gold patches, credentials, prompts, or responses.

\section{Methodology}
The exploratory protocol planned baseline and ArchGuard conditions for three qualified tasks. Four
PySnooper executions reached the provider. Both black-4 conditions are pre-provider infrastructure
attrition following checkout validation failure; they are neither provider outcomes nor repair
failures. Analysis is descriptive because only two valid pairs and zero discordant successful pairs
exist. No significance test is meaningful.

\section{Results}
Four of six executions completed and two attrited (33.3\%). Baseline and ArchGuard each achieved 0/2;
all four completed responses were invalid patches. Thus no effectiveness difference was observed, but
equivalence is not established. Baseline used 7,821 input and 1,037 output tokens at USD 0.0351075.
ArchGuard used 27,610 input and 1,700 output tokens at USD 0.094525. Total usage was 38,168 tokens and
USD 0.1296325. The higher observed ArchGuard cost cannot be generalized beyond these executions.

\begin{figure}[t]\centering\includegraphics[width=\columnwidth]{../figures/execution-outcomes.png}
\caption{Exploratory accounting: four invalid patches and two infrastructure attritions.}\end{figure}

\section{Failure Analysis and Discussion}
Evidence establishes patch-policy rejection but not a unique cause. Formatting, patch grammar,
policy, task context, integration, and model behavior remain possible mechanisms. Strict validation
improves safety and evidential clarity while potentially reducing apparent yield. Black-4 instead
reveals checkout robustness limitations. RQ1 receives engineering integration support; RQ2 identifies
patch validity and infrastructure as limitations; RQ3 records zero successes and condition-specific
usage. Repair effectiveness remains unresolved.

\section{Threats to Validity}
Construct validity depends on a strict, project-specific success definition. Internal validity is
limited by unknown invalid-patch causes and environment dependence. External validity is limited by
two pairs from one project. Conclusion validity precludes inference or superiority. Infrastructure
validity is affected by 33.3\% attrition. Results are model-, provider-, time-, and configuration-
dependent. Static analysis may miss runtime relations; profiles require human judgment; sandboxing
cannot eliminate all risk; and canonical evidence cannot reproduce stochastic provider behavior.

\section{Ethics and Reproducibility}
Negative outcomes and costs are retained. Only verified citations are used. Human authors remain
responsible for interpretation, disclosure, and submission. The final result fingerprint is
\texttt{761d9853f6ecfb4dafdaad81de781c4d49108b8b82912aeec509e173c7d56c70}.

\section{Conclusion and Future Work}
ArchGuard-APR demonstrates integration of architecture constraints with an auditable APR pipeline, but
does not demonstrate repair advantage. Future work should improve admissible patch production,
diagnostics, checkout robustness, and evaluation breadth before any confirmatory study.

\section*{Data Availability}Canonical source-free data are under \texttt{research/results/phase-13}.
\section*{Artifact Availability}Code, builders, tables, figures, and checks accompany the manuscript.
\section*{Author Contributions}HUMAN\_CONFIRMATION\_REQUIRED.
\section*{Funding}HUMAN\_CONFIRMATION\_REQUIRED.
\section*{Conflict of Interest}HUMAN\_CONFIRMATION\_REQUIRED.
\bibliographystyle{IEEEtran}\bibliography{archguard-apr-en}
\end{document}
